\documentclass[11pt]{article}

% ---------------------------------------------------------------------------
%  reharness: компилятор рассуждений для AI-воркфлоу
%  Научное описание рантайма и компилятора (русская версия, v0.25.0).
%  Сборка (нужен Unicode-движок):  tectonic reharness.ru.tex
%                              или: xelatex / lualatex reharness.ru.tex
% ---------------------------------------------------------------------------

\usepackage[a4paper,margin=1in]{geometry}
\usepackage{amsmath,amssymb,amsthm}
\usepackage{mathtools}
\usepackage{enumitem}
\usepackage{booktabs}
\usepackage{xcolor}

% ---- Unicode / кириллица (XeTeX/LuaTeX через tectonic) --------------------
% fontspec + babel (НЕ polyglossia: она строго требует OpenType-тег "Cyrl",
% которого нет в системных .ttc на macOS). XeTeX рисует кириллицу из шрифта
% напрямую; babel даёт русские переносы. На другой системе замените шрифты на
% любые установленные с поддержкой кириллицы.
\usepackage{fontspec}
\setmainfont{Times New Roman}
\setsansfont{Arial}
\setmonofont{Menlo}[Scale=0.9]
\usepackage[english,russian]{babel}

\usepackage{hyperref}
\hypersetup{colorlinks=true,linkcolor=blue!55!black,citecolor=green!45!black,urlcolor=blue!55!black}

% ---- окружения для теорем -------------------------------------------------
\theoremstyle{plain}
\newtheorem{theorem}{Теорема}[section]
\newtheorem{proposition}[theorem]{Утверждение}
\newtheorem{lemma}[theorem]{Лемма}
\newtheorem{corollary}[theorem]{Следствие}
\theoremstyle{definition}
\newtheorem{definition}[theorem]{Определение}
\newtheorem{principle}[theorem]{Принцип}
\newtheorem{invariant}[theorem]{Инвариант}
\theoremstyle{remark}
\newtheorem{remark}[theorem]{Замечание}

% ---- сокращения -----------------------------------------------------------
% Times New Roman не имеет small-caps начертания; набираем имя обычным шрифтом.
\newcommand{\reharness}{\textsf{reharness}}
\newcommand{\code}[1]{\texttt{#1}}
\newcommand{\state}[1]{\textsf{#1}}
\newcommand{\transition}{\ensuremath{\delta}}
\newcommand{\Nat}{\mathbb{N}}
\newcommand{\given}{\,\vert\,}
\newcommand{\powerset}{\mathcal{P}}
\newcommand{\meet}{\sqcap}
\newcommand{\joinop}{\sqcup}
\newcommand{\bigmeet}{\bigcap}

\title{\reharness: компилятор рассуждений для AI-воркфлоу\\[4pt]
\large Формальное описание рантайма и компилятора (v0.25.0)}
\author{проект reharness}
\date{\today}

\begin{document}
\maketitle

\begin{abstract}
\reharness{} компилирует запрос на естественном языке в исполнимый AI-воркфлоу на
основе конечных автоматов. Это \emph{компилятор рассуждений}: недетерминированный
интеллект больших языковых моделей (LLM) ограничен одноразовой фазой компиляции и
чётко очерченными \emph{агентскими листьями} в остальном полностью детерминированного
скелета исполнения. Система состоит из двух слоёв, оба --- конечные автоматы (КА).
\emph{Рантайм} исполняет скомпилированный КА и является именованным ограничением
стандартной, хорошо изученной модели: \emph{детерминированного иерархического
Мур-автомата-преобразователя с семантикой выполнения до завершения (run-to-completion)}.
\emph{Компилятор} превращает запрос в единый декларативный артефакт --- граф с
поведенческими контрактами на каждом узле --- и понижает его до исполнимого TypeScript
через детерминированные проходы по графу, проверяемые статическим анализатором,
построенным на двух классических движках (достижимость через транзитивное замыкание и
монотонный фреймворк потоков данных Кама--Ульмана). Поток данных между стадиями
\emph{выводится} из топологии графа через полиэдральную поэкземплярную модель, а не
декларируется. Мы приводим формальные модели, формулируем и доказываем гарантии
детерминизма и завершимости, и представляем устройство компилятора как небольшой набор
принципов, из которых выводится вся механическая машинерия. Мы соотносим систему с
предшествующими работами о компиляции LLM (DSPy, LLM+P, «Compiled AI»),
со statechart'ами, монотонным анализом потоков данных и полиэдральной моделью.
\end{abstract}

\tableofcontents

% ===========================================================================
\section{Введение}
% ===========================================================================

Агентские LLM-фреймворки обычно вызывают модель на \emph{каждом} шаге воркфлоу, смешивая
две различные задачи: \emph{решать, что делать} (где рассуждение полезно) и
\emph{исполнять чётко заданную процедуру} (где оно не нужно). Это смешение --- источник
трёх хорошо задокументированных проблем: недетерминизма (выходы меняются от прогона к
прогону даже при температуре $0$), стоимости (токены растут с объёмом исполнения) и
неаудируемости (нет фиксированного артефакта для инспекции).

\reharness{} занимает противоположную позицию, разделяемую небольшим корпусом недавних
работ о трактовке «LLM как компилятор, а не интерпретатор»~\cite{dspy,llmp,compiledai}:
интеллект тратится \emph{один раз}, на этапе компиляции, чтобы получить
детерминированный артефакт, который затем исполняется обычным программным обеспечением.
Там, где интеллект на этапе исполнения действительно неустраним --- суждение над
неструктурированным входом --- он изолируется в \emph{агентский лист}: один вызов LLM,
являющийся одним состоянием детерминированного КА.

\paragraph{Два слоя.} Архитектура --- это два конечных автомата:
\begin{itemize}[nosep]
  \item \textbf{рантайм} (\S\ref{sec:runtime}, \code{src/runtime/}): исполняет
        скомпилированный КА --- состояния суть действия agent/code/$\dots$, переходы
        типизированы и снабжены охранными условиями;
  \item \textbf{компилятор} (\S\ref{sec:compiler}, \code{src/compiler/}):
        мета-конвейер \code{/generate} превращает запрос в \code{skeleton.xml} и понижает
        его до исполнимого TypeScript-конвейера; анализатор в
        \code{src/compiler/analysis/} статически его проверяет.
\end{itemize}
Сам компилятор --- это \reharness{}-конвейер, исполняемый тем же рантаймом; система
самохостируема.

\paragraph{Инженерная философия.} Оба слоя суть \emph{именованные ограничения
стандартных моделей}, а не произвольные конструкции. Рантайм --- ограничение
иерархического автомата Харела/UML; анализатор --- семейство экземпляров двух
стандартных движков; модель потоков данных --- полиэдральное пространство итераций.
Ведущий инженерный принцип: \emph{упрощай, пока не сломается детерминизм}: кодовая база
неоднократно \emph{удаляла} машинерию (декларируемый поток данных, циклы ревью,
сквозные блоки спецификаций), когда шов обходился дороже проверки, которую он покупал.
Этот документ фиксирует получившуюся теорию.

% ===========================================================================
\section{Рантайм}\label{sec:runtime}
% ===========================================================================

Рантайм (\code{src/runtime/fsm.ts}) исполняет скомпилированный конвейер. Формально это
\textbf{детерминированный иерархический Мур-автомат-преобразователь, управляемый шагами
выполнения до завершения (RTC)}: детерминированное конечное управление со стеком
рекурсии для составных состояний (ограниченная магазинная структура), где действие
каждого состояния --- самодостаточное вычисление, испускающее ровно одно событие. Это
намеренное ограничение иерархического автомата UML/Харела (HSM) --- та же топология
управления за вычетом машинерии, чей недетерминизм мы отвергаем
(\S\ref{sec:restrictions}).

\subsection{Синтаксис конвейера}

\begin{definition}[Конвейер]\label{def:pipeline}
Конвейер --- это кортеж $P = (Q, q_0, F, \tau, \delta_{\mathrm{on}}, G, \kappa)$, где
\begin{itemize}[nosep]
  \item $Q$ --- конечное множество \emph{состояний}, $q_0 \in Q$ --- \emph{начальное};
  \item $F \subseteq Q$ --- \emph{финальные} состояния, каждое помечено \code{success} или \code{error};
  \item $\tau : Q \to \{\state{agent}, \state{code}, \state{interactive}, \state{switch},
        \state{set}, \state{parallel}, \state{loop}, \state{call}, \state{wait}, \state{final}\}$
        задаёт \emph{тип} состояния;
  \item $\delta_{\mathrm{on}} : Q \times \Sigma \rightharpoonup (Q \times G)^{*}$
        сопоставляет состоянию и \emph{событию} $e \in \Sigma$ упорядоченный список
        охраняемых целей;
  \item $G$ --- множество \emph{охранных условий} (булевых выражений над хранилищем данных, \S\ref{sec:guards});
  \item $\kappa$ задаёт составным состояниям их структурные параметры (тела ветвей/шагов,
        цель join, верхнюю границу итераций и т.д.).
\end{itemize}
\end{definition}

Состояния делятся на \emph{активные} (\state{agent}, \state{code}, \state{interactive}),
чья семантика --- \emph{действие}; \emph{управляющие} (\state{switch}, \state{set});
\emph{составные} (\state{parallel}, \state{loop}, \state{call}); состояние внешнего
сигнала \state{wait}; и \state{final}.

\subsection{Мур-семантика и шаг выполнения до завершения}

\begin{definition}[Мур-действие]
Каждое активное состояние $q$ несёт \emph{действие входа} $a_q$, которое выполняется до
завершения и \emph{испускает ровно одно событие} $e \in \Sigma$, равное \emph{итогу
собственного вычисления}. Испускаемое событие не является функцией входного символа,
прочитанного в середине шага (это был бы автомат Мили); это результат $a_q$
(свойство Мура~\cite{moore}).
\end{definition}

Конфигурация --- пара $(q, \sigma)$, где $q$ --- активное состояние, а $\sigma$ ---
\emph{хранилище данных} (отображение ключ/значение \code{ctx.data} в памяти, плюс
read-only \code{config} и пространство рабочих директорий стадий; \S\ref{sec:dataflow}).
Рантайм выполняет по одному шагу RTC за раз.

\begin{definition}[Шаг RTC]\label{def:rtc}
Шаг выполнения-до-завершения из конфигурации $(q,\sigma)$ при активном $\tau(q)$:
\begin{enumerate}[nosep]
  \item выполнить $a_q$ до завершения, получив событие $e$ и обновлённое хранилище $\sigma'$;
  \item вычислить следующее состояние $q' = \transition(q, e, \sigma')$ (\S\ref{sec:delta});
  \item следующий шаг начинается только после полного завершения текущего.
\end{enumerate}
Нет ни очереди событий, ни вытеснения, ни чередования: в шаге существует ровно одно
событие, так что требование RTC «нет внутреннего параллелизма внутри шага» выполняется
тривиально. Состояния \state{wait} --- единственный источник событий, синхронизируемых
извне.
\end{definition}

\subsection{Функция переходов}\label{sec:delta}

\begin{definition}[Функция переходов]\label{def:delta}
Для состояния $q$, испущенного события $e$ и хранилища $\sigma$ пусть
$\delta_{\mathrm{on}}(q,e) = \langle (t_1,g_1), \dots, (t_k,g_k)\rangle$ --- упорядоченный
список охраняемых целей. Тогда
\[
  \transition(q,e,\sigma) \;=\; t_{j}, \qquad
  j = \min\{\, i \given g_i(\sigma) = \text{истина} \,\},
\]
т.е.\ цель \emph{первого} охранного условия, истинного в порядке документа. Если
$\delta_{\mathrm{on}}(q,e)$ не определено или ни одно $g_i$ не истинно, рантайм
\emph{громко падает}: останавливается с диагностикой, а не молча зависает.
\end{definition}

\begin{definition}[Тотальность через завершение]\label{def:total}
Конвейер \emph{тотален}, если для всякой достижимой конфигурации $(q,\sigma)$ и всякого
события $e$, которое может испустить $a_q$, $\transition(q,e,\sigma)$ определено и
какое-то охранное условие истинно. Кодген и линтер обеспечивают структурные условия,
достаточные для тотальности (у каждого испускаемого события есть переход; у состояний
\state{switch} есть ветвь по умолчанию; существует финал \state{error} для
синтезируемого ребра \code{ERROR}); где они не могут, $\transition$ громко падает
(Опр.~\ref{def:delta}), так что рантайм никогда не зависает молча. Мы трактуем сам «громкий
сбой» как делающий $\transition$ тотальной в $Q \cup \{\bot\}$, где $\bot$ ---
наблюдаемый аварийный останов.
\end{definition}

\subsection{Язык охранных условий}\label{sec:guards}

Охранные условия $g \in G$ берутся из намеренно крошечного, свободного от побочных
эффектов подъязыка выражений, чтобы они были статически анализируемы и сами не вносили
недетерминизма или произвольных вычислений. Грамматика допускает: корневые идентификаторы
\code{config}, \code{data}, \code{retries} с доступом к членам; операторы сравнения и
булевы ${==}\ {!=}\ {<}\ {\le}\ {>}\ {\ge}\ \&\&\ \vert\vert\ !$ и арифметические
${+}\ {-}\ {*}\ {/}\ \%$; литералы строк/чисел/булевых/\code{null}; и литералы массивов.
Она \emph{отвергает} вызовы функций, присваивание, тернарный оператор,
инкремент/декремент, побитовые операторы и регулярные выражения. Охрана повтора имеет
особый вид $\code{retries:}k < N$ и читает ограниченный счётчик повторов $k$. Компилятор
понижает охрану в чистое тело стрелочной функции TypeScript (корневые идентификаторы
получают префикс контекстного объекта); рантайм никогда не вычисляет охрану как строку.
Запрет вызовов --- также свойство безопасности: текст охраны происходит из авторского XML
от LLM и встраивается дословно в сгенерированный код, так что допущенный вызов был бы
поверхностью исполнения кода (\S\ref{sec:limitations}).

\subsection{Иерархия: составные состояния как подвычисления RTC}

Составные состояния исполняются как вложенные подвычисления RTC со своим кадром рекурсии
и возвращают родителю единственное завершение; переход родителя трактует весь составной
узел как одно Мур-действие.

\paragraph{\state{parallel} (fork--join).}
Состояние \state{parallel} вычисляет выражение \code{over} в массив
$[x_0,\dots,x_{m-1}]$, исполняет тело \emph{ветви} по разу на элемент (с ограничением на
параллелизм), и после завершения всех ветвей переходит к своему \code{join}. Агентские
ветви исполняются с \emph{настоящим параллелизмом ОС-процессов} (каждая порождает
отдельный процесс \code{pi}); чисто-кодовые ветви исполняются кооперативно в едином цикле
событий. Join читает результаты по ветвям $\code{data.branches} = [\{i, x_i,
\mathit{dir}_i, \mathit{ok}_i\}]$ как барьер.

\begin{invariant}[Ветвь не пишет в общее хранилище]\label{inv:branch}
Ветвь \state{parallel} не должна писать в общее хранилище \code{ctx.data}. Ветви
обмениваются данными только через свои изолированные выходные директории;
межветвевые данные читает join после барьера.
\end{invariant}

\paragraph{\state{loop} (ограниченная итерация).}
\state{loop} исполняет упорядоченный список состояний-\emph{шагов} по разу на итерацию.
Он несёт обязательную жёсткую границу \code{max} $\in \Nat$ и необязательный предикат
раннего выхода \code{exit}. Рантайм поддерживает две различные величины пространства
итераций (\S\ref{sec:loopcounters}).

\paragraph{\state{call} (подавтомат).}
\state{call} вызывает другой скомпилированный конвейер как подконвейер и переходит по его
терминальному статусу (\code{success}/\code{error}).

\subsection{Детерминизм}

\begin{theorem}[Детерминизм]\label{thm:determinism}
Зафиксируем конвейер $P$ и выходы всех агентских листьев. Тогда наблюдаемый результат
исполнения $P$ --- последовательность посещённых состояний, финальное хранилище и
множество записанных файлов рабочего пространства --- является детерминированной функцией
начального хранилища, с точностью до коммутативного слияния выходов ветвей \state{parallel}.
\end{theorem}

\begin{proof}[Набросок доказательства]
Покажем, что каждый источник выбора разрешается детерминированно.
(i) Внутри шага $a_q$ выполняется до завершения и испускает единственное событие
(Опр.~\ref{def:rtc}); единственный недетерминизм внутри $a_q$ --- агентский лист, который
зафиксирован по предположению. (ii) $\transition$ --- функция, а не отношение: на пару
$(q,e)$ срабатывает не более одного перехода, потому что охранные условия сканируются в
порядке документа и побеждает \emph{первое} истинное (Опр.~\ref{def:delta}); нет
неоднозначности «высшего приоритета» и нет одновременно разрешённой гонки. (iii) Составные
состояния возвращают единственное завершение, так что родитель видит одно Мур-действие.
(iv) Для \state{parallel}: ветви не могут писать в общее \code{ctx.data}
(Инв.~\ref{inv:branch}), а join --- барьер, поэтому единственная зависимость от порядка ---
порядок слияния выходов ветвей; но поскольку каждая ветвь пишет в непересекающуюся,
индексируемую по номеру выходную директорию (\S\ref{sec:dataflow}), слияние коммутативно и
наблюдаемый результат не зависит от порядка завершения. Композиция (i)--(iv) по
последовательности шагов RTC даёт детерминированный преобразователь с точностью до этого
коммутативного слияния.
\end{proof}

\begin{remark}
Теорема~\ref{thm:determinism} --- центральный контракт устройства: КА есть
\emph{детерминированный скелет} вокруг недетерминированных агентских листьев.
«Compiled AI» в смысле Trooskens et al.\ доводит число листьев до нуля (никаких LLM на
этапе исполнения); \reharness{}-конвейер сохраняет листья там, где рассуждение
неустранимо, но держит \emph{обвязку} вокруг них детерминированной. Полностью
детерминированный случай (ноль агентских состояний) --- достижимая точка устройства, а не
отдельный режим (\S\ref{sec:amortization}).
\end{remark}

\subsection{Завершимость}

\begin{invariant}[Циклы ограничены]\label{inv:max}
Каждое состояние \state{loop} объявляет жёсткую границу $\code{max} \in \Nat$,
$\code{max} \ge 1$ (проверяется линтером, \S\ref{sec:lint}). Предикат раннего выхода
\code{exit} --- необязательное уточнение, никогда не заменяющее границу.
\end{invariant}

\begin{theorem}[Завершимость]\label{thm:termination}
Всякое исполнение конвейера $P$, в котором каждый агентский лист завершается,
останавливается за конечное число шагов RTC.
\end{theorem}

\begin{proof}[Набросок доказательства]
Определим вполне-обоснованный ранг на конфигурациях. \state{loop} вносит счётчик,
ограниченный сверху своим \code{max} (Инв.~\ref{inv:max}); он строго растёт на каждой
итерации и принуждает к выходу при \code{max}, так что цикл делает не более \code{max}
итераций. \state{parallel} делает $m = |\code{over}|$ исполнений ветвей, $m$ конечно.
\state{call} рекурсирует в строго меньший подконвейер (граф вызовов ацикличен по
построению). Циклы верхнего уровня в графе управления допустимы лишь под охраной счётчика
\code{retries}, ограниченного своим \code{retries-max} (ребро повтора --- единственный
способ вернуться в раннее состояние, а счётчик монотонен). Таким образом, всякое обратное
ребро ограничено статическим счётчиком, и лексикографическое произведение этих счётчиков
есть вполне-обоснованная мера, строго убывающая вдоль всякого шага, который иначе мог бы
повториться. Конечное число неповторяющихся шагов плюс вполне-обоснованная мера на
повторяющихся даёт конечность, при условии что каждое $a_q$ (в частности, каждый
агентский лист) завершается.
\end{proof}

\subsection{Счётчики цикла: координата против мощности}\label{sec:loopcounters}

Тонкий, но поучительный момент. Цикл выставляет в хранилище данных две
\emph{категориально различные} величины, и их смешение --- латентная ошибка
«сдвиг на единицу».

\begin{definition}[Координата и мощность цикла]
Для цикла с итерациями, индексируемыми $0,1,\dots$:
\begin{itemize}[nosep]
  \item \code{data.iteration} --- \textbf{координата}: $0$-индексированный номер
        \emph{текущей} итерации. Это компонента вектора экземпляра
        (\S\ref{sec:dataflow}), действительная \emph{внутри} шага. Предикат раннего
        выхода \code{exit} видит её как индекс только что завершившейся итерации.
  \item \code{data.iterations} --- \textbf{мощность}: число выполненных итераций. Это
        возвращаемое значение цикла, доступное \emph{после} цикла.
\end{itemize}
\end{definition}

Различие зеркалит разницу между индексом массива $i$ и его длиной $|A|$. «Сколько раундов
прошло» --- это мощность \code{data.iterations}; её нельзя вычислять как
$\code{data.iteration}+1$, а \code{data.iteration} нельзя читать после цикла. Рантайм
восстанавливает в общий скаляр после каждой итерации координату (а не мощность), чтобы
вложенный шаг цикла, перезаписавший её, не испортил представление объемлющего цикла.

\subsection{Ограничения относительно полного UML HSM}\label{sec:restrictions}

Рантайм \reharness{} намеренно опускает из полной модели statechart'ов:
\begin{itemize}[nosep]
  \item \textbf{ортогональные регионы с широковещанием событий} --- заменены на
        fork--join \state{parallel} \emph{без} межветвевой сигнализации. Широковещание ---
        главный источник недетерминизма statechart'ов и взрыва пространства состояний;
  \item \textbf{Мили-действия на переходах и цепочки entry/exit-действий} --- всё
        поведение живёт в Мур-действиях состояний; переходы --- чистое управление;
  \item \textbf{внутренние/локальные переходы, глубокая история, гонки событий
        завершения} --- опущены.
\end{itemize}
Остаётся ровно анализируемое ядро: граф управления, стек рекурсии и действия на
состояниях, испускающие одно событие. Именно это делает Теоремы~\ref{thm:determinism} и
\ref{thm:termination} верными по построению.

% ===========================================================================
\section{Модель потоков данных и рабочего пространства}\label{sec:dataflow}
% ===========================================================================

Центральное проектное решение: \emph{поток данных между стадиями выводится из топологии
графа, а не декларируется}. Нет атрибута \code{produces}/\code{consumes}, нет синтаксиса
путей, нет ручной проводки. Сперва обоснуем это исторически, затем дадим полиэдральную
модель, которая это реализует.

\subsection{Почему выводится, а не декларируется}

Два более ранних устройства декларировали поток данных --- сначала файловые артефакты с
путями \code{artifact=}/\code{requires=}, затем слой \code{produces}/\code{consumes} ---
и оба \emph{переносили} шов дрейфа, а не устраняли его: строки путей расходятся с
реальностью, а дрейф «декларация против прозы» рецидивирует на стыке слияния
параллельных стадий. Осознание:

\begin{principle}[Ребро есть контракт данных]\label{prin:edge}
Ребро графа $A \to B$ уже \emph{есть} утверждение, что $B$ может читать выход $A$. Поэтому
автор не декларирует о проводке ничего; компилятор выводит её из рёбер.
\end{principle}

\subsection{Три канала}

Данные движутся по трём каналам:
\begin{enumerate}[nosep]
  \item \textbf{\code{config}} --- внешний CLI-интерфейс (\S\ref{sec:inputs}), read-only,
        \emph{декларируется} (у него нет внутреннего производителя, из которого можно вывести);
  \item \textbf{\code{ctx.data}} --- скаляры в памяти, пишутся только состояниями
        \state{code} и \state{set}, читаются охранными условиями / \code{over} /
        \code{exit} / \code{model-expr};
  \item \textbf{рабочие директории стадий} --- файлы; каждая стадия получает свою выходную
        директорию и читает директории своих видимых производителей.
\end{enumerate}

\subsection{Полиэдральная поэкземплярная модель}

Модель рабочего пространства --- это \emph{полиэдральное пространство итераций}
(поэкземплярный анализ потоков данных по массивам Фотрие~\cite{feautrier}): стадия
исполняется по разу на точку своего пространства итераций, а чтение видит экземпляры
производителя на префиксе общей области видимости, собранные по осям, которые потребитель
покинул.

\begin{definition}[Объемлющая область / вектор экземпляра]\label{def:scope}
Для состояния $n$ его \emph{объемлющая область} $\mathrm{scope}(n) = \langle c_1,
\dots, c_d \rangle$ --- цепочка составных состояний, объемлющих его, от внешнего к
внутреннему (\state{parallel} вносит ось ветви, \state{loop} --- ось итерации). Один
\emph{экземпляр} $n$ адресуется \emph{вектором экземпляра} $\vec{\imath} = (i_1, \dots,
i_d) \in \Nat^d$, по одному индексу на ось. Его выходная директория есть
\[
  \mathrm{dir}(n, \vec{\imath}) \;=\; \langle\mathit{runDir}\rangle/\code{work}/n/i_1/i_2/\cdots/i_d .
\]
Стадия верхнего уровня имеет $d = 0$ и единственную директорию.
\end{definition}

\begin{definition}[Видимые производители]\label{def:visible}
Производители, видимые $n$, --- это его \emph{производители-предки}: состояния типа
\state{agent}/\state{code}/\state{interactive}, обратно достижимые из $n$ через
ролевое отношение преемников (\S\ref{sec:engines}). Каждый производитель $p$
классифицируется по мощности выхода относительно $n$. Пусть $\ell =
|\mathrm{lcp}(\mathrm{scope}(n),\mathrm{scope}(p))|$ --- длина наибольшего общего
префикса двух областей. Тогда
\[
  p \in
  \begin{cases}
    \mathit{single}(n) & \text{если } |\mathrm{scope}(p)| = \ell, \\[2pt]
    \mathit{list}(n)   & \text{иначе.}
  \end{cases}
\]
\end{definition}

Интуитивно: $p$ \emph{одиночный} тогда и только тогда, когда каждое составное состояние,
объемлющее $p$, объемлет и $n$ (экземпляр $p$ фиксирован объемлющим контекстом $n$ ---
линейные цепочки, чтения в одной области, чтения с переносом по циклу). $p$ ---
\emph{список} тогда и только тогда, когда $n$ покинул хотя бы одно составное состояние,
объемлющее $p$; тогда $n$ видит \emph{коллекцию} экземпляров $p$ по покинутым осям: по
директории на ветвь (покинутый \state{parallel} --- это \emph{map}) и/или на итерацию
(покинутый \state{loop} --- это \emph{scan} / история).

\begin{proposition}[Один закон для map, scan и вложенности]
Определение~\ref{def:visible} даёт корректную видимость и мощность единообразно для:
линейных цепочек (включая пропуск через дальних предков), fork-out/in \state{parallel}
(map), межитерационной истории \state{loop} (scan) и чтений с переносом по циклу, а также
вложенных составных состояний произвольной глубины (вектор экземпляра растёт на одну ось
на каждое составное состояние).
\end{proposition}

Во время исполнения три аксессора выставляют это стадии: \code{c.out()} --- её
собственная директория экземпляра; \code{c.dir(p)} --- единственный экземпляр
производителя из общей области (с учётом переноса по циклу: последний существующий
экземпляр $\le$ текущего индекса); \code{c.dirs(p)} --- коллекция по покинутым осям.
Запись производителя, записанный \code{data.branches[i].dir} и чтение потребителя --- все
выводятся из $(\,\text{стадия}, \vec{\imath}\,)$, поэтому не могут разойтись.

\begin{remark}[Граница]
Топология выводит \emph{проводку} и \emph{мощность} (какие экземпляры видимы), а не
\emph{форму} (схему файла --- поведенческое свойство, проверяемое проходом \state{polish})
и не \emph{избирательность} (какое подмножество стадии следует читать, управляемое прозой
контракта). Условные производители на альтернативных ветвях \state{switch} предлагаются
пермиссивно (нет гарантии «произведено на каждом пути» для файлов; только \code{ctx.data}
несёт проверку обязательной записи \S\ref{sec:defassign}).
\end{remark}

% ===========================================================================
\section{Статический анализатор}\label{sec:compiler-analysis}
% ===========================================================================

Всякая статическая проверка \reharness{} --- это \emph{экземпляр} одного из двух
переиспользуемых движков (\code{src/compiler/analysis/framework.ts}). Новые проверки
добавляются как экземпляры; кодовая база не пишет вручную самодельные неподвижные точки.

\subsection{Два движка}\label{sec:engines}

\begin{definition}[Движок достижимости]
$\mathrm{reachableFrom}(S, \mathrm{next})$ вычисляет транзитивное замыкание множества $S$
относительно отношения преемников $\mathrm{next}$ обходом в ширину с рабочим списком.
Решётка тривиальна (множество достижимого).
\end{definition}

\begin{definition}[Движок монотонного анализа потоков данных]\label{def:kamullman}
$\mathrm{solveMonotoneSets}$ --- это \textbf{монотонный фреймворк потоков данных
Кама--Ульмана}~\cite{kamullman}, в русле унифицированной решёточной формулировки
глобальной оптимизации программ Килдалла~\cite{kildall}, специализированный к решётке
$L = (\powerset(K), \subseteq)$ подмножеств конечного множества ключей $K$ с
gen-only прямой передаточной функцией
\[
  \mathrm{OUT}[n] = \mathrm{gen}(n) \cup \mathrm{IN}[n], \qquad
  \mathrm{IN}[n] = \!\!\bigmeet_{p \in \mathrm{pred}(n)}\!\! \mathrm{OUT}[p],
\]
итерируемой до неподвижной точки, где $\meet$ --- пересечение множеств (анализ
\textbf{must}) или объединение множеств (анализ \textbf{may}). Достигающие определения,
живые переменные и определённое присваивание --- классические
экземпляры~\cite{dragon,nielson}.
\end{definition}

\begin{proposition}[Сходимость]\label{prop:converge}
$\mathrm{solveMonotoneSets}$ завершается в наименьшей неподвижной точке. Передаточные
функции монотонны, а решётка $\powerset(K)$ имеет конечную высоту $|K|$; следовательно, по
теореме Кама--Ульмана хаотическая итерация сходится.
\end{proposition}

\subsection{Проверки как экземпляры}

\begin{table}[h]
\centering
\small
\begin{tabular}{@{}lll@{}}
\toprule
\textbf{Проверка} & \textbf{Экземпляр движка} & \textbf{Модуль} \\
\midrule
достижимость от начального  & прямая достижимость от $q_0$              & \code{semantic.ts} \\
достижимость до финала      & обратная достижимость от $F$ ( co-reach.) & \code{semantic.ts} \\
определённое присваивание   & прямой \emph{must} (meet $=$ пересечение)& \code{dataflow.ts} \\
\code{visibleProducers}     & обратная may-дост.\ $+$ мощность          & \code{graph.ts} \\
поток config                & множество чтений $\subseteq$ объявленные  & \code{dataflow.ts} \\
грамматика / корректность   & «система типов» формата                   & \code{lint.ts} \\
покрытие контрактами        & у каждого активного есть \code{<contract>}& \code{semantic.ts} \\
\bottomrule
\end{tabular}
\caption{Каждая статическая проверка --- именованный экземпляр движка из \S\ref{sec:engines}.}
\label{tab:checks}
\end{table}

\subsection{Достижимость и тупики}

Достижимость от начального --- это $\mathrm{reachableFrom}(\{q_0\}, \mathrm{succ})$; всякое
состояние не в результате недостижимо. Достижимость до финала ---
$\mathrm{reachableFrom}(F, \mathrm{succ}^{-1})$ по обратным рёбрам (ко-достижимость);
всякое состояние не в результате есть тупик. Отношение преемников \emph{ролевое}: ветвь
\state{parallel} или шаг \state{loop} маршрутизируется через join родительского
составного состояния, а не через свои рёбра \code{<on>}.

\subsection{Определённое присваивание над \code{ctx.data}}\label{sec:defassign}

\begin{definition}[Определённое присваивание]
Чтение \code{data.k} в состоянии $n$ \emph{безопасно}, если \code{k} записывается на
\emph{каждом} пути, ведущем в $n$. Это учебниковая прямая \emph{must}-задача потоков
данных: экземпляр Опр.~\ref{def:kamullman} с $\meet = {}$пересечение и $\mathrm{gen}(n)$,
равным ключам, которые $n$ определённо пишет. Пусть $\mathrm{IN}[n]$ --- неподвижная
точка. Чтение \code{data.k} в $n$ помечается, если $\code{k} \notin \mathrm{IN}[n]$.
\end{definition}

\begin{proposition}[Корректность must-проверки]
Если проверка не помечает чтение \code{data.k} в $n$, то \code{k} записывается на каждом
пути управления от $q_0$ до $n$; следовательно, \code{data.k} определена всякий раз, когда
$n$ исполняется. Эквивалентно: проверка помечает чтение лишь тогда, когда доказуемо
существует путь без записывающего.
\end{proposition}

Уточнение различает \emph{чтения-на-входе} (объявленные чтения; выражение \code{over}
состояния \state{parallel}, вычисляемое на входе; выражения значений состояния
\state{set}) и \emph{чтения-после} (охранные условия на исходящих переходах самого
состояния, вычисляемые \emph{после} того, как действие входа выполнилось). Чтения-на-входе
проверяются против $\mathrm{IN}[n]$; чтения-после --- против $\mathrm{IN}[n] \cup
\mathrm{gen}(n)$, поскольку собственные записи узла доступны его собственным исходящим
охранным условиям. Это устраняет ложное срабатывание на частом шаблоне «состояние
\state{code} пишет \code{data.x}, затем маршрутизируется по \code{expr:data.x}».

% ===========================================================================
\section{Внешний интерфейс}\label{sec:inputs}
% ===========================================================================

Единственное, что о данных \emph{декларируется}, --- это CLI-интерфейс, потому что он
внешний: он приходит от пользователя, а не из какого-либо ребра графа, так что нет
производителя, из которого его вывести.

\begin{principle}[Декларируй внешнее, выводи внутреннее]\label{prin:law}
Внутренний поток данных выводится из рёбер графа (Принцип~\ref{prin:edge}). Внешний
интерфейс (CLI-входы) декларируется --- в блоке \code{<inputs>} на конвейер. Это
неустранимая сигнатура конвейера, а не возврат к декларируемому потоку данных.
\end{principle}

Каждый \code{<arg>} объявляет поле \code{config.X}; кодген генерирует парсер аргументов из
объявления (позиционный или флаг, приведение типов среди
\code{string}/\code{list}/\code{number}/\code{bool}, значения по умолчанию,
обязательность). Два поля всегда предоставляются и никогда не объявляются:
\code{config.target} (рабочая директория) и \code{config.input} (позиционные аргументы,
склеенные вместе).

\begin{definition}[Проверка потока config]
Пусть $\mathrm{Decl}$ --- объявленные входы $\cup\, \{\code{target}, \code{input}\}$, а
$\mathrm{Used}$ --- множество полей \code{config.X}, читаемых где-либо (выражения скелета:
охранные условия, \code{over}, \code{exit}, \code{model-expr}, значения \state{set}; и
сгенерированный код: \code{c.config.X}). Проверка требует $\mathrm{Used} \subseteq
\mathrm{Decl}$.
\end{definition}

\begin{corollary}[Корректность CLI по построению]
Если конвейер проходит проверку потока config, то каждое поле \code{config.X}, читаемое
конвейером, разбирается из командной строки сгенерированным парсером. Следовательно,
«компилируется» влечёт «проводка CLI корректна».
\end{corollary}

% ===========================================================================
\section{Конвейер компилятора}\label{sec:compiler}
% ===========================================================================

Компилятор (\code{/generate}) --- самохостируемый \reharness{}-конвейер: это КА,
исполняемый рантаймом из \S\ref{sec:runtime}. Его стадии:
\begin{center}\small
\state{maybe\_research} $\to$ \state{research} $\to$ \state{prd} $\to$
\state{review\_prd} \emph{(approval)} $\to$ \state{design} $\to$ \state{construct}
$\to$ \state{fill\_prompts} $\to$ \state{check\_dataflow} $\to$ \state{polish}
$\to$ \state{verify} $\to$ \state{done}.
\end{center}

\subsection{Принцип минимализма и единый артефакт}

\begin{principle}[LLM авторствует минимальный артефакт]\label{prin:lean}
LLM пишет ровно один артефакт --- \code{skeleton.xml}: \emph{граф} (состояния $+$
типизированные переходы) плюс поведенческий \code{<contract>} на каждом узле (CDATA). Всё
механическое (парсер аргументов, проводка между стадиями, каркас команды и библиотеки) ---
детерминированный проход по графу. Нет \code{produces}/\code{consumes}, нет
\code{<spec>}, нет ручных путей.
\end{principle}

\subsection{Человек одобряет намерение, а не граф}

\begin{principle}[PRD-first]\label{prin:prd}
Единственная человеческая контрольная точка --- \state{review\_prd}: человек правит и
одобряет \emph{PRD} (документ требований к продукту: цель, поведение, критерии приёмки,
границы) --- формулировку \emph{намерения}. Человек никогда не одобряет граф КА, который
есть механическая реализация, принадлежащая компилятору. Намерение --- то, что человек
может оценить; реализация --- то, что может проверить анализатор.
\end{principle}

PRD --- также точка \emph{унификации} источников входа: \code{/generate} отображает
$\text{промпт} + \text{рисёрч} \to \text{PRD}$, а планируемый фронтенд \code{compile}
отображает $\text{логи сессии} + \text{рисёрч} \to \text{PRD}$; всё после PRD
переиспользуется без изменений.

\subsection{Понижение и верификация}

\state{construct} разбирает и статически проверяет скелет (lint $+$ semantic,
Табл.~\ref{tab:checks}) и эмитит команду, заглушки библиотеки code-состояний и заглушки
промптов агентов. \state{fill\_prompts} превращает каждый \code{<contract>} либо в промпт
агента (для состояний \state{agent}), либо в реализацию кода (для состояний \state{code}),
параллельно. \state{check\_dataflow} запускает анализатор потоков данных (определённое
присваивание $+$ поток config). \state{verify} компилирует сгенерированный TypeScript
(\code{tsc --noEmit}) и загружает граф КА каждой команды. Обратите внимание, что
верификация \emph{статична}: компилятор никогда не исполняет сгенерированный конвейер;
поведенческая проверка против прогонов --- удел отдельного, будущего контура обратной
связи (\S\ref{sec:limitations}).

\subsection{Коррекция: один корректор в горячем контексте}\label{sec:lint}

\begin{principle}[Polish, а не цикл ревью]\label{prin:polish}
Коррекция --- это один агент \state{polish}, который проверяет весь сгенерированный
конвейер против PRD и чинит \emph{листья} (промпты агентов, реализации кода) в одном
горячем контексте --- заменяя более старый цикл ревью$\to$починка$\to$переревью.
\state{redesign} --- редкая крайняя мера для изменений топологии (когда \state{polish} не
может починить проблему в листьях и эскалирует); он остаётся верным одобренному PRD ---
меняет \emph{как}, никогда не \emph{что}.
\end{principle}

\begin{principle}[Решения принадлежат состояниям или агентам]\label{prin:decisions}
Решения о маршрутизации принадлежат состояниям КА (\state{switch}, охраняемые переходы);
решения-суждения принадлежат агенту. Ни то, ни другое не принадлежит императивному клею
между стадиями. Рантайм не разбирает, не классифицирует и не шлюзует между ходами агента.
\end{principle}

% ===========================================================================
\section{Амортизация рассуждений}\label{sec:amortization}
% ===========================================================================

Состояние \state{code} --- это \emph{замороженное рассуждение}: написано один раз на этапе
компиляции, затем исполняется при нулевой стоимости в токенах навсегда. Амортизация --- это
политика компилятора реализовывать стадию как \state{code}, а не \state{agent}, всякий раз,
когда рассуждение на самом деле не является неустранимым. Поскольку выполняется инвариант
«узел целиком \state{code} или целиком \state{agent}» (встроенный вызов LLM внутри
code-состояния отвергается линтером), амортизация --- это \emph{полное} понижение
$\state{agent} \to \state{code}$, и её граница объективна.

\begin{definition}[Понижаемость]\label{def:demote}
Стадия $n$ \emph{понижаема} до \state{code} тогда и только тогда, когда выполнены оба
условия:
\begin{enumerate}[label=(\alph*),nosep]
  \item \textbf{структурный вход}: каждый производитель, читаемый $n$ (его предки по
        Опр.~\ref{def:visible}), есть состояние \state{code}/\state{set}, так что байты
        несут схему, авторскую от компилятора. \emph{Это условие статически разрешимо} из
        графа.
  \item \textbf{механический контракт}: контракт --- тотальная вычислимая функция над этой
        структурой (переформатировать, отфильтровать, отсортировать, дедуплицировать,
        подсчитать, арифметика, извлечь поле, проверить схему). Это \emph{не} суждение
        (резюмировать, оценить, классифицировать неоднозначный текст, кластеризовать «тот
        же ли это вопрос?», назначить серьёзность). \emph{Это условие статически
        неразрешимо} --- по теореме Райса~\cite{rice} семантическое свойство «эта
        спецификация --- чистая функция над структурой» неразрешимо --- так что это
        суждение, которое по Принципу~\ref{prin:decisions} принадлежит агенту-дизайнеру, а
        не статическому правилу.
\end{enumerate}
\end{definition}

\begin{remark}[Асимметрия и опасность пере-понижения]
Условие (a) необходимо, но не достаточно: можно совершать суждение над JSON. Понижение при
ложном (b) --- худшая ошибка, эмпирически задокументированный режим отказа, в котором
регулярное выражение над счётом-фактурой извлекает имя поставщика с точностью $\sim$20\%
там, где LLM достигает $\sim$80\%: вход выглядел разбираемым, но задача была семантической.
Поэтому политика консервативна: при сомнении в (b) оставляем стадию \state{agent}. Правило
реализовано внутри промпта \state{design} (оно не стоит дополнительного прохода и
дополнительных токенов); суждение о (b) выносится один раз, на этапе авторства.
\end{remark}

Оба предельных случая достижимы: чисто генеративный линейный конвейер (например, рисёрч
$\to$ план $\to$ рендер) не допускает понижения сверх финального ввода-вывода, потому что
цепочка повсюду питается выходом агентов; чисто механический конвейер (например,
проверить-затем-склеить два файла) понижается до \emph{нуля} агентских состояний ---
полностью детерминированный воркфлоу с нулём токенов, что есть в точности предел
«Compiled AI», достигнутый со стороны agent-first.

% ===========================================================================
\section{Связь с предшествующими работами}
% ===========================================================================

\paragraph{LLM-как-компилятор.} DSPy~\cite{dspy} компилирует декларативные LLM-конвейеры в
оптимизированные промпты/веса; LLM+P~\cite{llmp} переводит естественный язык в формальную
спецификацию планирования (PDDL~\cite{pddl}), делегируемую классическому решателю;
«Compiled AI»~\cite{compiledai} генерирует детерминированный код на фазе компиляции и
исполняет его при нуле вызовов модели в рантайме, с четырёхстадийным конвейером
генерации-и-валидации и библиотекой компонентов «шаблон/модуль/блок-промпта». \reharness{}
разделяет позицию «компилировать один раз», но отличается тремя способами: (i) он
компилирует \emph{оркестрацию} (анализируемый КА с выводимым потоком данных), а не эмитит
бизнес-код в фиксированные шаблоны; (ii) он сохраняет неустранимое рассуждение как
анализируемые агентские листья, а не запрещает любой вывод в рантайме (его «Code
Factory»/bounded-agentic запасной путь --- в \reharness{} основной и единственный режим,
управляемый Опр.~\ref{def:demote}); и (iii) человек одобряет \emph{намерение} (PRD), а не
конфигурацию.

\paragraph{Конечные автоматы.} Рантайм --- именованное ограничение
\emph{Statechart'ов} Харела~\cite{harel} и семантики выполнения-до-завершения
\emph{UML}~\cite{uml,scxml,samek}, отбрасывающее ортогональные регионы и Мили-действия на
рёбрах (\S\ref{sec:restrictions}); состояния несут Мур-действия~\cite{moore}. Конструкция
\state{parallel} --- модель задач fork--join, а не параллелизм ортогональных регионов.

\paragraph{Анализ программ.} Анализатор --- монотонный фреймворк потоков данных
Кама--Ульмана~\cite{kamullman,kildall} (достигающие определения, живые переменные и
определённое присваивание --- его классические экземпляры~\cite{dragon,nielson}) плюс
достижимость через транзитивное замыкание; модель потоков данных/рабочего пространства ---
полиэдральное пространство итераций (поэкземплярный анализ потоков данных по массивам
Фотрие~\cite{feautrier}). Неразрешимость семантического условия понижаемости
(Опр.~\ref{def:demote}) следует из теоремы Райса~\cite{rice}.

% ===========================================================================
\section{Ограничения и границы}\label{sec:limitations}
% ===========================================================================

\begin{itemize}[nosep]
  \item \textbf{Компиляция статична по построению.} Компилятор никогда не запускает
        сгенерированный конвейер; \state{verify} проверяет типы и загрузку графа, а не
        поведение. Поведенческая корректность (например, правило консенсуса, корректное по
        форме, но вакуумное; краевой случай в code-состоянии) наблюдается только из
        прогонов и является уделом отдельного контура обратной связи (\emph{evolve}),
        соединённого с компилятором внешним циклом $\text{compile} \to \text{run} \to
        \text{evolve}$, а не исполнением внутри \code{/generate}.
  \item \textbf{Источник намерения изменчив.} Короткий промпт сильно опирается на
        (недетерминированную) стадию веб-рисёрча, поэтому PRD --- а значит, и топология ---
        меняется от прогона к прогону. Это изменчивость намерения, а не дрейф реализации;
        контрольная точка одобрения PRD человеком (Принцип~\ref{prin:prd}) --- якорь против
        неё.
  \item \textbf{Условные файловые производители} на альтернативных ветвях \state{switch}
        не имеют гарантии обязательного производства (только \code{ctx.data} несёт проверку
        обязательной записи); выходы подконвейеров \state{call} не моделируются как
        производители.
  \item \textbf{Форма и избирательность вне области топологии.} Топология выводит проводку
        и мощность, а не схемы файлов или то, какое подмножество читать; это поведенческие
        свойства, проверяемые \state{polish} или заявленные в контрактах.
\end{itemize}

% ===========================================================================
\section{Заключение}
% ===========================================================================

\reharness{} --- компилятор рассуждений, чьи два слоя суть именованные ограничения
стандартных моделей: детерминированный иерархический Мур-автомат-преобразователь для
исполнения (детерминированный и завершимый по построению,
Теоремы~\ref{thm:determinism} и \ref{thm:termination}), и минималистичный, PRD-first
компилятор, чья всякая статическая проверка --- экземпляр фреймворка Кама--Ульмана или
достижимости через транзитивное замыкание, и чей поток данных между стадиями выводится из
полиэдрального пространства итераций, а не декларируется. Объединяющий закон ---
\emph{выводи внутреннее, декларируй внешнее} --- и правило амортизации --- \emph{понижай в
код тогда и только тогда, когда вход структурен, а контракт механичен} --- позволяют
одному и тому же конвейеру охватить весь спектр от чисто генеративного (все агенты) до
чисто детерминированного (ноль агентов, предел Compiled-AI). Дисциплина устройства ---
держать этот скелет детерминированным и анализируемым и тратить недетерминированный
интеллект только там, где он неустраним.

% ===========================================================================
\begin{thebibliography}{99}
% ===========================================================================

\bibitem{compiledai} G.~Trooskens, A.~Karlsberg, A.~Sharma, L.~De Brouwer,
M.~Van Puyvelde, M.~Young, J.~Thickstun, G.~Alterovitz, and W.~A.~De Brouwer.
\emph{Compiled AI: Deterministic Code Generation for LLM-Based Workflow
Automation}. arXiv:2604.05150, 2026.

\bibitem{dspy} O.~Khattab et al.
\emph{DSPy: Compiling Declarative Language Model Calls into Self-Improving
Pipelines}. arXiv:2310.03714, 2023.

\bibitem{llmp} B.~Liu, Y.~Jiang, X.~Zhang, Q.~Liu, S.~Zhang, J.~Biswas, and
P.~Stone. \emph{LLM+P: Empowering Large Language Models with Optimal Planning
Proficiency}. arXiv:2304.11477, 2023.

\bibitem{pddl} D.~McDermott et al.
\emph{PDDL --- The Planning Domain Definition Language}.
Tech.\ Report CVC TR-98-003, Yale Center for Computational Vision and Control, 1998.

\bibitem{harel} D.~Harel.
\emph{Statecharts: A Visual Formalism for Complex Systems}.
Science of Computer Programming, 8(3):231--274, 1987.

\bibitem{uml} Object Management Group.
\emph{OMG Unified Modeling Language (OMG UML), Version 2.5.1}.
OMG, 2017. (State Machines; run-to-completion semantics.)

\bibitem{scxml} World Wide Web Consortium.
\emph{State Chart XML (SCXML): State Machine Notation for Control Abstraction}.
W3C Recommendation, 2015.

\bibitem{samek} M.~Samek.
\emph{Practical UML Statecharts in C/C++: Event-Driven Programming for Embedded
Systems}. 2nd ed., Newnes, 2008.

\bibitem{moore} E.~F.~Moore.
\emph{Gedanken-Experiments on Sequential Machines}.
In \emph{Automata Studies}, Annals of Mathematics Studies 34, pp.~129--153,
Princeton University Press, 1956.

\bibitem{kamullman} J.~B.~Kam and J.~D.~Ullman.
\emph{Monotone Data Flow Analysis Frameworks}.
Acta Informatica, 7(3):305--317, 1977.

\bibitem{kildall} G.~A.~Kildall.
\emph{A Unified Approach to Global Program Optimization}.
In \emph{Proc.\ 1st ACM SIGACT-SIGPLAN Symposium on Principles of Programming
Languages (POPL)}, pp.~194--206, 1973.

\bibitem{dragon} A.~V.~Aho, M.~S.~Lam, R.~Sethi, and J.~D.~Ullman.
\emph{Compilers: Principles, Techniques, and Tools}. 2nd ed., Addison-Wesley, 2006.
(Гл.~9, анализ потоков данных.)

\bibitem{nielson} F.~Nielson, H.~R.~Nielson, and C.~Hankin.
\emph{Principles of Program Analysis}. Springer, 1999.

\bibitem{feautrier} P.~Feautrier.
\emph{Dataflow Analysis of Array and Scalar References}.
International Journal of Parallel Programming, 20(1):23--53, 1991.

\bibitem{rice} H.~G.~Rice.
\emph{Classes of Recursively Enumerable Sets and Their Decision Problems}.
Transactions of the American Mathematical Society, 74(2):358--366, 1953.

\end{thebibliography}

\appendix
\section{Обозначения}
\begin{center}\small
\begin{tabular}{@{}ll@{}}
\toprule
Символ & Значение \\
\midrule
$Q, q_0, F$ & состояния, начальное состояние, финальные состояния \\
$\Sigma$ & алфавит событий (\code{DONE}, \code{ERROR}, \code{PASS}, $\dots$) \\
$\transition$ & функция переходов (Опр.~\ref{def:delta}) \\
$\sigma$ & хранилище данных (\code{ctx.data} $+$ \code{config} $+$ рабочее пространство) \\
$\mathrm{scope}(n)$ & цепочка объемлющих составных состояний $n$ (Опр.~\ref{def:scope}) \\
$\vec{\imath}$ & вектор экземпляра, по одному индексу на ось \\
$\mathit{single}(n), \mathit{list}(n)$ & видимые производители по мощности (Опр.~\ref{def:visible}) \\
$L = (\powerset(K), \subseteq)$ & решётка потоков данных (Опр.~\ref{def:kamullman}) \\
$\meet$ & meet ($\cap$ для must, $\cup$ для may) \\
\bottomrule
\end{tabular}
\end{center}

\end{document}
